Put
\[
 v_-:=\min_{d\in\mathcal D}\min\{V_{E,d},V_{O,d}\}>0,
 \qquad
 Z_m:=\max_{d\in\mathcal D}\max\bigl\{
 |\widehat V_{E,d}-V_{E,d}|,
 |\widehat V_{O,d}-V_{O,d}|\bigr\}.
\]
The strict inequality follows from the nondegeneracy assumption and finiteness of \(\mathcal D\).  For \(x,y>0\), define
\[
 F_d(x,y):=\left(\sqrt{c_{E,d}x}+\sqrt{c_{O,d}y}\right)^2.
\]
The budget-profile theorem gives
\(F_d(V_{E,d},V_{O,d})=\mathcal V_d^*(P)\), and the pilot construction gives
\(F_d(\widehat V_{E,d},\widehat V_{O,d})=
\widehat{\mathcal V}_{d,m}^*\).  On the event \(Z_m<v_-/2\), the arguments of every \(F_d\) lie in a fixed compact subset of \((0,\infty)^2\).  The gradients of the finitely many maps \(F_d\) are uniformly bounded there.  The mean-value theorem and the uniform pilot variance-rate assumption therefore give
\[
 \varepsilon_m\le C Z_m,
 \qquad
 \varepsilon_m=O_P(r_{V,m})=o_P(1).                                      \tag{1}
\]

Fix any \(d^*\in D^*(P)\).  The empirical optimality of \(\widehat d_m\) implies, pathwise,
\[
\begin{aligned}
0
&\le \mathcal V_{\widehat d_m}^*(P)-\mathcal V_{d^*}^*(P)\\
&\le
 \left|\mathcal V_{\widehat d_m}^*(P)
       -\widehat{\mathcal V}_{\widehat d_m,m}^*\right|
 +\left\{\widehat{\mathcal V}_{\widehat d_m,m}^*
       -\widehat{\mathcal V}_{d^*,m}^*\right\}
 +\left|\widehat{\mathcal V}_{d^*,m}^*
       -\mathcal V_{d^*}^*(P)\right|\\
&\le 2\varepsilon_m.
\end{aligned}                                                              \tag{2}
\]
If \(D^*(P)\ne\mathcal D\), finiteness of \(\mathcal D\) implies that every design outside \(D^*(P)\) has loss at least \(\gamma(P)>0\).  Consequently, (2) gives
\[
 \{\widehat d_m\notin D^*(P)\}
 \subseteq \{2\varepsilon_m\ge\gamma(P)\}.                                \tag{3}
\]
Equations (1) and (3) prove the probability assertion.  If \(D^*(P)=\mathcal D\), the left-hand event in (3) is empty and the right-hand event is empty under the convention \(\gamma(P)=+\infty\).

For the allocation, write
\[
 A_d(x,y):=
 \frac{\sqrt{c_{E,d}x}}
 {\sqrt{c_{E,d}x}+\sqrt{c_{O,d}y}}.
\]
The elementary inequality
\[
 |(x\vee\underline v_m)-V|
 \le |x-V|+\underline v_m
\]
and the mean-value theorem, again on a common positive compact rectangle and uniformly over the finite family, imply
\[
 \max_{d\in\mathcal D}
 |\widehat\alpha_{d,m}-\alpha_d^*(P)|
 =O_P(r_{V,m}+\underline v_m).                                              \tag{4}
\]
The second-order allocation-loss lemma now yields the sharper uniform bound
\[
 \max_{d\in\mathcal D}\left|
 B\left\{
 \frac{V_{E,d}}{\lfloor\widehat\alpha_{d,m}B/c_{E,d}\rfloor}
 +\frac{V_{O,d}}{\lfloor(1-\widehat\alpha_{d,m})B/c_{O,d}\rfloor}
 \right\}-\mathcal V_d^*(P)\right|
 =O_P\!\left((r_{V,m}+\underline v_m)^2+B^{-1}\right).                    \tag{5}
\]
Since \(r_{V,m}+\underline v_m=o(1)\), (5) implies the theorem's stated weaker order
\(O_P(r_{V,m}+\underline v_m+B^{-1})=o_P(1)\).  The deterministic tie rule chooses a single empirical minimizer, whereas (2)--(3) correctly compare that choice with the entire population optimizer set.